\section{Paged allocation and shared prefixes}
\label{sec:C04:paging}
PagedAttention separates logical token positions from physically contiguous storage. Fixed-size blocks hold keys and values, and a block table maps logical positions to physical blocks. This avoids reserving one large contiguous buffer for the maximum possible length of every request. It also permits sharing of completed prefix blocks across related sequences, with copy-on-write behavior when a shared mutable region would otherwise diverge. The original vLLM paper attributes throughput gains to improved memory use and associated batching opportunities under its evaluated workloads. Those are author-reported results for that implementation and its baselines~\cite{kwon2023pagedattention}.

Let a block hold $B$ tokens. A sequence of length $T$ occupies $B\lceil T/B\rceil$ token slots. Its internal slack is less than $B$ slots; across $N$ sequences the total slack is less than $NB$, before considering sharing or metadata. This bound explains one advantage over reserving each request's maximum length. It does not imply zero overhead. Smaller blocks reduce terminal slack but increase block-table size and indirection frequency. Larger blocks can improve locality while wasting more space near sequence ends. The best block size depends on the length distribution, attention kernel, and sharing patterns.

Prefix sharing is valid only for an identical computation prefix. Identical visible text is insufficient if tokenization, positions, masks, model weights, adapters, or relevant execution settings differ. Even an identical token sequence can yield different keys and values when an earlier hidden input differs. A robust cache key therefore binds the complete computation identity required by the model implementation, together with tenant isolation rules. A semantic embedding match can suggest useful retrieval; it cannot justify substituting exact KV tensors. Confusing these mechanisms can create silent correctness failures rather than merely a lower cache hit rate.

Editing a prefix invalidates more than the edited tokens in ordinary causal attention. Representations at later positions depend on the earlier prefix through successive layers. Reusing their old keys and values after a deletion or replacement generally evaluates a different computational graph from a fresh forward pass. Safe reuse requires a proved invariant, a model architecture that supports the transformation, or recomputation of affected state. Context-editing systems therefore need explicit interaction with the cache manager. The compression service cannot be treated as a purely textual component whose output is transparently compatible with every retained tensor.

An agent often branches on uncertain plans or candidate patches. Shared prefixes can reduce physical memory duplication, but every divergent continuation still consumes computation and new blocks. A search controller that increases the branch count can exhaust capacity even while the cache manager is efficient. Admission control should reserve an envelope for likely continuation or support preemption with declared semantics. Recompute-on-resume and host offload are different policies with different latency and bandwidth costs. Neither can be ranked without the session length distribution, tool-pause distribution, and deadline requirements.

